\documentclass{zhvt-classic}

% 使用本地源流明体（GenRyuMin2 TW）按文件路径加载，不依赖系统字体名
% 版框（内框／界栏／书口侧栏／鱼尾）改用与版心印章同一印泥棕色：
%   印章图片 yshs.png 实色 = RGB(135,68,52) = HTML 874434
\usepackage{xcolor}
\usepackage{tikz}
% sealbrown 双注册：xcolor 供 TikZ/pgf，L3 供 cls 版框
\definecolor{sealbrown}{HTML}{874434}
\ExplSyntaxOn
\color_set:nnn { sealbrown } { HTML } { 874434 }
\ExplSyntaxOff

\zhvtset{
  grid color       = sealbrown,
  main font        = GenRyuMin2TW-SB,
  main font path   = C:/Users/Administrator/texmf/fonts/opentype/public/GenRyuMin2TW/,
jiazhu font      = GenRyuMin2TW-M,
  jiazhu font path = C:/Users/Administrator/texmf/fonts/opentype/public/GenRyuMin2TW/,
  % 页码框容纳 4 字：正常页「九十七」+前置空 1 格=4 行恰好；
  % 100+ 页「一百零一」共 5 行超出，末字「一/二/三」被框裁掉（2026-10-04 实测），故放宽到 5
  shukou page number chars = 5,
  % 书口书名框高从默认 6 字减到 4 字，把省出的 2 字高拨给 chapter 框，
  % 使 6 字书名+卷名放进 9 字高的 chapter 框后不再压线（2026-10-05）
  shukou title chars = 4,
}

\title{史记}[陈涉世家]
\author{汉·司马迁}
\maker{闲叟漫记·重製}


% ===== 索隐述赞样式（2026-11-08 四次定稿 · 愚大人亲手改定参数）=====
% ⚠ 竖排前提：自定义字体家族必须带 RawFeature={vertical:+vert}；
%   simhei.ttf 加 +vhal 会整体加载失败静默回退，只写 +vert；simfang 同理。
\newCJKfontfamily{\shuzanfont}[Path=C:/Windows/Fonts/,RawFeature={vertical:+vert}]{simhei.ttf}
\newCJKfontfamily{\shuzanbody}[Path=C:/Windows/Fonts/,RawFeature={vertical:+vert}]{simfang.ttf}
% 题头：152x30pt 圆角 6pt 棕底白字；baseline (0,2) 使整框落入界行格子不压版框线；
% 白字节点 x 不补偿、y +8.2pt（像素级墨迹实测 dx/dy≈0，L/R 留白约 20pt）。
\newcommand{\shuzanhead}[1]{%
  \begin{tikzpicture}[baseline={(0,2)}]
    \fill[sealbrown, rounded corners=6pt]
      (-76pt,-15pt) rectangle (76pt,15pt);
    \node[text=white, inner sep=0pt, anchor=center] at (0pt,8.2pt)
      {\shuzanfont\fontsize{23}{23}\selectfont #1};
  \end{tikzpicture}%
}
% ★ 修正书口卷名位置：chapter 框起始位置改回 cls 默认值，避免全局影响。
%   只在每个卷首页通过 \juan 命令里的空命令微调该页位置（安全，不会跨页污染）。
\ExplSyntaxOn
\AtBeginDocument{
  \dim_set:Nn \l__zhvt_shukou_chapter_title_start_dim
    { \l__zhvt_shukou_title_height_dim + \l__zhvt_font_size_dim * 2 }
  \dim_set:Nn \l__zhvt_shukou_page_number_start_dim
    { \l__zhvt_shukou_chapter_title_start_dim + \l__zhvt_shukou_chapter_title_height_dim }
}
\ExplSyntaxOff

\begin{document}

% ★★ 版心（书口）排布覆写：把书名放到鱼尾之下（2026-10-04）
%
%   cls 默认版心自上而下：书名（title 静态框，在鱼尾之上）→ 鱼尾 → 篇名（chapter 动态框）→ 页码
%   本模板改为：          鱼尾（位置不动）→ 书名 → 篇名 → 页码分割线 → 页码
%
%   手法：把 cls 放在 title 静态框里的书名，改放进 chapter 动态框（鱼尾之下）；
%         title 静态框留空。书名取 \l__zhvt_title_str，即 \title{...} 的第一参数，
%         所以换书时只改 \title，不必在此重复写书名。
%
%   ⚠ 三条必须遵守的写法（2026-10-04 实测）：
%     ① 汉字必须用 \hbox{...} 包住。flowfram 会把框内容装进一个宽 14.4mm（= 一个 line_sep
%        = 书口栏宽）的纵向 minipage 并按该宽度折行，而整页靠 PDF /Rotate 90 显示、
%        框内文字本身不变向 —— 不包 \hbox 时「围炉夜话」4 字会被逐字折成 4 行，
%        在版面上就成了横排 4 字并横向溢出版框。
%     ② 竖排下**一个 \hbox 就是一整列**。书名与篇名要同处一列，必须放进同一个 \hbox；
%        分成两个 \hbox 时，后一个会被当成新的一列、甩到隔壁栏去
%        （实测「序」跑到书口栏外侧 14.4mm 的正文列上）。中间的 \quad 在框内是横向空白，
%        显示出来正好是书名与篇名之间的纵向间隔。
%     ③ ★ 内容盒要写成 **与框等宽的定宽盒** `\hbox to \l__zhvt_shukou_chapter_title_height_dim{…\hfil}`。
%        实测：chapter 框是把内容盒沿纵向**居中**摆放的，内容宽度一变，书名就跟着漂：
%        自然宽度的 \hbox（书名 38.4mm）落在 x=157.2mm，加篇名后（57.6mm）落在 x=152.4mm，
%        目录页（无篇名）与正文页的书名就差半格。做成等宽盒后，无论有没有篇名，
%        内容固定从框顶起排，实测两页均落在 x=142.8mm，不再漂。
\ExplSyntaxOn
% 把书口 chapter 框（书名+篇名）起始位置下移 1 字高，避免 4 字卷名顶到鱼尾
% ⚠ cls 的 \__zhvt_calculate_dimensions: 在 \AtBeginDocument 里执行，会重算该 dim；
%   必须把覆写挂到 begin document 之后，否则会被二次覆盖（2026-10-05 实测）。
% (旧方案已移到导言区，直接修改 chapter 框起始位置)
% 卷标题命令 #1=书口篇名(鱼尾) #2=正文卷名（含卷序，卷名内空格用 ~）
%
% ★★ 目录页两项修订（2026-10-05，愚大人指出）★★
%   ① 页码起点统一对齐：
%      原状 —— 目录项写成 \par\hspace{4em}卷名\hspace{4em}页码，页码起点 = 卷名实宽 + 4em。
%      卷名实宽随字数变：「卷一~梁惠王上」6 字 +~、「卷七~离娄上」5 字 +~，
%      差 1 个字位 → 卷七~卷十 4 条的页码整体比其余 10 条高 27.21pt（一个字），看着上下不齐。
%      改法 —— 卷名补全角空格（宽 = 1 字 = 27.215pt）补足到 **7 字符位**再写进目录，
%      14 条页码起点遂全部落在同一条竖线上（起点对齐；末字随页码字数自然参差）。
%      ⚠ 只对目录里的副本 \l_tmpc_tl 补，正文卷名仍用 \l_tmpa_tl 原样 —— 正文像素不受影响。
%   ② 卷名整体上移半字位：前导 \hspace 由 4em 改 3.5em。
%      实测原状目录块在版框内偏下（天头留白 139.6pt / 地脚留白 110.7pt），
%      上移 0.5em = 13.6pt 后上下基本均分，视觉居中。
\NewDocumentCommand { \juan } { mm } {
	\clearpage
	\tl_set:Nn \l__zhvt_chapter_title_tl { #1 }
	\tl_set:Nn \l_tmpa_tl { #2 }
	\tl_set_eq:NN \l_tmpc_tl \l_tmpa_tl
	\int_step_inline:nn { 7 - \str_count:n { #2 } }
	  { \tl_put_right:Nn \l_tmpc_tl { 　 } }
	\tl_set:Nx \l_tmpb_tl { \zhnum { page } }
	\tl_put_left:Nn \l_tmpb_tl { \hspace { 4em } }
	\tl_put_left:Nn \l_tmpb_tl { \l_tmpc_tl }
	\tl_put_left:Nn \l_tmpb_tl { \par \hspace { 3.5em } }
	\addtocontents { toc } { \l_tmpb_tl }
	\skip_horizontal:n { \l__zhvt_font_size_dim * 4 }
	\l_tmpa_tl
	\par
	\thispagestyle { empty }
}
\cs_set_protected:Npn \__zhvt_set_shukou_content: {
	% ★ 鱼尾上方留白（title 静态框）嵌入闲叟钤印：
	%   整页 /Rotate 90，故此处把图用 \rotatebox{90} 转正为竖章；
	%   尺寸卡在版心条幅内（书内约 7mm 宽 × 33.8mm 高，留余量不压鱼尾）。
	\setstaticcontents * { title } {
		\vbox_to_ht:nn { \l__zhvt_baseline_skip_dim } {
			\vskip 14.2pt
			\hbox_to_wd:nn { \l__zhvt_shukou_title_height_dim } {
				\hfil
				\rotatebox [ origin = c ] { 90 } { \includegraphics [ width = 7mm , height = 33.8mm ] { yshs.png } }
				\hfil
			}
		}
	}
% ★ 版心（书口）统一改法（2026-11-16 愚大人定）：18pt＋框内居中＋书口文字同版框色
%   ① 原 27.2pt 时 7 字长篇名（如「廉颇蔺相如列传」）末字溢出 chapter 框、压住页码；
%      改 \fontsize{18pt}{18pt}，前后 \hfil 令字串在框内居中（不随篇名字数漂移）。
%   ② \vbox{\vskip 8.2pt} 定量下移，抵消缩字号后栏宽方向 8.2pt 上偏（像素实测 y 中心 543.3 = 原版）。
%      ⚠ \raisebox / \vspace* 在 flowfram 框内无效；换字号须同步换补偿：16pt→10.1、19pt→7.4。
%   ③ chapter 与 page~number 均置 \color{sealbrown}（HTML 874434），与版框／界行同色；正文仍黑。
%   ⚠ 页码字位、钤印、鱼尾、界行、正文一律未动；五部合订本同此口径。
	\setdynamiccontents * { chapter }
	{ \color{sealbrown} \vbox { \vskip 8.2pt \hbox to \l__zhvt_shukou_chapter_title_height_dim
		{ \hfil \fontsize{18pt}{18pt}\selectfont \l__zhvt_title_str \quad \l__zhvt_chapter_title_tl \hfil } } }
	\setdynamiccontents * { page~number } { \color{sealbrown} \quad \zhnum { page } }
}
\__zhvt_set_shukou_content:
\ExplSyntaxOff

\maketitle

% 单篇，无多卷目录，不渲染目录页

\juan{陈涉世家}{卷四十八~陈涉世家第十八}
\hspace*{2em}陈胜者。阳城人也。字涉。吴广者。阳夏人也。字叔。陈涉少时。尝与人佣耕。辍耕之垄上。怅恨久之。曰。﹁苟富贵。无相忘。﹂庸者笑而应曰。﹁若为庸耕。何富贵也。﹂陈涉太息曰。﹁嗟乎。燕雀安知鸿鹄之志哉。﹂\\
\hspace*{2em}二世元年七月。发闾左谪戍渔阳。九百人屯大泽乡。陈胜。吴广皆次当行。为屯长。会天大雨。道不通。度已失期。失期。法皆斩。陈胜。吴广乃谋曰。﹁今亡亦死。举大计亦死。等死。死国可乎。﹂陈胜曰。﹁天下苦秦久矣。吾闻二世少子也。不当立。当立者乃公子扶苏。扶苏以数谏故。上使外将兵。今或闻无罪。二世杀之。百姓多闻其贤。未知其死也。项燕为楚将。数有功。爱士卒。楚人怜之。或以为死。或以为亡。今诚以吾众诈自称公子扶苏。项燕。为天下唱。宜多应者。﹂吴广以为然。乃行卜。卜者知其指意。曰。﹁足下事皆成。有功。然足下卜之鬼乎。﹂陈胜。吴广喜。念鬼。曰。﹁此教我先威众耳。﹂乃丹书帛曰﹁陈胜王﹂。置人所罾鱼腹中。卒买鱼烹食。得鱼腹中书。固以怪之矣。又间令吴广之次所旁丛祠中。夜篝火。狐鸣呼曰﹁大楚兴。陈胜王﹂。卒皆夜惊恐。旦日。卒中往往语。皆指目陈胜。\\
\hspace*{2em}吴广素爱人。士卒多为用者。将尉醉。广故数言欲亡。忿恚尉。令辱之。以激怒其众。尉果笞广。尉剑挺。广起。夺而杀尉。陈胜佐之。并杀两尉。召令徒属曰。﹁公等遇雨。皆已失期。失期当斩。藉弟令毋斩。而戍死者固十六七。且壮士不死即已。死即举大名耳。王侯将相宁有种乎。﹂徒属皆曰。﹁敬受命。﹂乃诈称公子扶苏。项燕。从民欲也。袒右。称大楚。为坛而盟。祭以尉首。陈胜自立为将军。吴广为都尉。攻大泽乡。收而攻蕲。蕲下。乃令符离人葛婴将兵徇蕲以东。攻铚。酂。苦。柘。谯皆下之。行收兵。比至陈。车六七百乘。骑千余。卒数万人。攻陈。陈守令皆不在。独守丞与战谯门中。弗胜。守丞死。乃入据陈。数日。号令召三老。豪杰与皆来会计事。三老。豪杰皆曰。﹁将军身被坚执锐。伐无道。诛暴秦。复立楚国之社稷。功宜为王。﹂陈涉乃立为王。号为张楚。\\
\hspace*{2em}当此时。诸郡县苦秦吏者。皆刑其长吏。杀之以应陈涉。乃以吴叔为假王。监诸将以西击荥阳。令陈人武臣。张耳。陈余徇赵地。令汝阴人邓宗徇九江郡。当此时。楚兵数千人为聚者。不可胜数。\\
\hspace*{2em}葛婴至东城。立襄强为楚王。婴后闻陈王已立。因杀襄强。还报。至陈。陈王诛杀葛婴。陈王令魏人周市北徇魏地。吴广围荥阳。李由为三川守。守荥阳。吴叔弗能下。陈王征国之豪杰与计。以上蔡人房君蔡赐为上柱国。\\
\hspace*{2em}周文。陈之贤人也。尝为项燕军视日。事春申君。自言习兵。陈王与之将军印。西击秦。行收兵至关。车千乘。卒数十万。至戏。军焉。秦令少府章邯免郦山徒。人奴产子生。悉发以击楚大军。尽败之。周文败。走出关。止次曹阳二三月。章邯追败之。复走次渑池十余日。章邯击。大破之。周文自刭。军遂不战。\\
\hspace*{2em}武臣到邯郸。自立为赵王。陈余为大将军。张耳。召骚为左右丞相。陈王怒。捕系武臣等家室。欲诛之。柱国曰。﹁秦未亡而诛赵王将相家属。此生一秦也。不如因而立之。﹂陈王乃遣使者贺赵。而徙系武臣等家属宫中。而封耳子张敖为成都君。趣赵兵亟入关。赵王将相相与谋曰。﹁王王赵。非楚意也。楚已诛秦。必加兵于赵。计莫如毋西兵。使使北徇燕地以自广也。赵南据大河。北有燕。代。楚虽胜秦。不敢制赵。若楚不胜秦。必重赵。赵乘秦之弊。可以得志于天下。﹂赵王以为然。因不西兵。而遣故上谷卒史韩广将兵北徇燕地。\\
\hspace*{2em}燕故贵人豪杰谓韩广曰。﹁楚已立王。赵又已立王。燕虽小。亦万乘之国也。愿将军立为燕王。﹂韩广曰。﹁广母在赵。不可。﹂燕人曰。﹁赵方西忧秦。南忧楚。其力不能禁我。且以楚之强。不敢害赵王将相之家。赵独安敢害将军之家。﹂韩广以为然。乃自立为燕王。居数月。赵奉燕王母及家属归之燕。\\
\hspace*{2em}当此之时。诸将之徇地者。不可胜数。周市北徇地至狄。狄人田儋杀狄令。自立为齐王。以齐反击周市。市军散。还至魏地。欲立魏后故甯陵君咎为魏王。时咎在陈王所。不得之魏。魏地已定。欲相与立周市为魏王。周市不肯。使者五反。陈王乃立甯陵君咎为魏王。遣之国。周市卒为相。\\
\hspace*{2em}将军田臧等相与谋曰。﹁周章军已破矣。秦兵旦暮至。我围荥阳城弗能下。秦军至。必大败。不如少遗兵。足以守荥阳。悉精兵迎秦军。今假王骄。不知兵权。不可与计。非诛之。事恐败。﹂因相与矫王令以诛吴叔。献其首于陈王。陈王使使赐田臧楚令尹印。使为上将。田臧乃使诸将李归等守荥阳城。自以精兵西迎秦军于敖仓。与战。田臧死。军破。章邯进兵击李归等荥阳下。破之。李归等死。\\
\hspace*{2em}阳城人邓说将兵居郏。章邯别将击破之。邓说军散走陈。铚人伍徐将兵居许。章邯击破之。伍徐军皆散走陈。陈王诛邓说。\\
\hspace*{2em}陈王初立时。陵人秦嘉。铚人董绁。符离人朱鸡石。取虑人郑布。徐人丁疾等皆特起。将兵围东海守庆于郯。陈王闻。乃使武平君畔为将军。监郯下军。秦嘉不受命。嘉自立为大司马。恶属武平君。告军吏曰。﹁武平君年少。不知兵事。勿听。﹂因矫以王命杀武平君畔。\\
\hspace*{2em}章邯已破伍徐。击陈。柱国房君死。章邯又进兵击陈西张贺军。陈王出监战。军破。张贺死。\\
\hspace*{2em}腊月。陈王之汝阴。还至下城父。其御庄贾杀以降秦。陈胜葬砀。谥曰隐王。\\
\hspace*{2em}陈王故涓人将军吕臣为仓头军。起新阳。攻陈下之。杀庄贾。复以陈为楚。\\
\hspace*{2em}初。陈王至陈。令铚人宋留将兵定南阳。入武关。留已徇南阳。闻陈王死。南阳复为秦。宋留不能入武关。乃东至新蔡。遇秦军。宋留以军降秦。秦传留至咸阳。车裂留以徇。\\
\hspace*{2em}秦嘉等闻陈王军破出走。乃立景驹为楚王。引兵之方与。欲击秦军定陶下。使公孙庆使齐王。欲与并力俱进。齐王曰。﹁闻陈王战败。不知其死生。楚安得不请而立王。﹂公孙庆曰。﹁齐不请楚而立王。楚何故请齐而立王。且楚首事。当令于天下。﹂田儋诛杀公孙庆。\\
\hspace*{2em}秦左右校复攻陈。下之。吕将军走。收兵复聚。鄱盗当阳君黥布之兵相收。复击秦左右校。破之青波。复以陈为楚。会项梁立怀王孙心为楚王。\\
\hspace*{2em}陈胜王凡六月。已为王。王陈。其故人尝与庸耕者闻之。之陈。扣宫门曰。﹁吾欲见涉。﹂宫门令欲缚之。自辩数。乃置。不肯为通。陈王出。遮道而呼涉。陈王闻之。乃召见。载与俱归。入宫。见殿屋帷帐。客曰。﹁伙颐。涉之为王沉沉者。﹂楚人谓多为伙。故天下传之。伙涉为王。由陈涉始。客出入愈益发舒。言陈王故情。或说陈王曰。﹁客愚无知。颛妄言。轻威。﹂陈王斩之。诸陈王故人皆自引去。由是无亲陈王者。陈王以朱房为中正。胡武为司过。主司群臣。诸将徇地。至。令之不是者。系而罪之。以苛察为忠。其所不善者。弗下吏。辄自治之。陈王信用之。诸将以其故不亲附。此其所以败也。陈胜虽已死。其所置遣侯王将相竟亡秦。由涉首事也。高祖时为陈涉置守冢三十家砀。至今血食。\\
\hspace*{2em}褚先生曰。地形险阻。所以为固也。兵革刑法。所以为治也。犹未足恃也。夫先王以仁义为本。而以固塞文法为枝叶。岂不然哉。吾闻贾生之称曰。\\
\hspace*{2em}﹁秦孝公据殽函之固。拥雍州之地。君臣固守。以窥周室。有席卷天下。包举宇内。囊括四海之意。并吞八荒之心。当是时也。商君佐之。内立法度。务耕织。修守战之备。外连衡而斗诸侯。于是秦人拱手而取西河之外。\\
\hspace*{2em}﹁孝公既没。惠文王。武王。昭王蒙故业。因遗策。南取汉中。西举巴蜀。东割膏腴之地。北收要害之郡。诸侯恐惧。会盟而谋弱秦。不爱珍器重宝肥饶之地。以致天下之士。合从缔交。相与为一。当此之时。齐有孟尝。赵有平原。楚有春申。魏有信陵。此四君者。皆明知而忠信。宽厚而爱人。尊贤而重士。约从连衡。兼韩。魏。燕。赵。宋。卫。中山之众。于是六国之士有宁越。徐尚。苏秦。杜赫之属为之谋。齐明。周冣。陈轸。邵滑。楼缓。翟景。苏厉。乐毅之徒通其意。吴起。孙膑。带他。儿良。王廖。田忌。廉颇。赵奢之伦制其兵。尝以什倍之地。百万之师。仰关而攻秦。秦人开关而延敌。九国之师遁逃而不敢进。秦无亡矢遗镞之费。而天下固已困矣。于是从散约败。争割地而赂秦。秦有余力而制其弊。追亡逐北。伏尸百万。流血漂橹。因利乘便。宰割天下。分裂山河。强国请服。弱国入朝。\\
\hspace*{2em}﹁施及孝文王。庄襄王。享国之日浅。国家无事。\\
\hspace*{2em}﹁及至始皇。奋六世之余烈。振长策而御宇内。吞二周而亡诸侯。履至尊而制六合。执敲朴以鞭笞天下。威振四海。南取百越之地。以为桂林。象郡。百越之君俯首系颈。委命下吏。\\
\hspace*{2em}﹁乃使蒙恬北筑长城而守藩篱。却匈奴七百余里。胡人不敢南下而牧马。士亦不敢贯弓而报怨。于是废先王之道。燔百家之言。以愚黔首。堕名城。杀豪俊。收天下之兵聚之咸阳。销锋镝。铸以为金人十二。以弱天下之民。然后践华为城。因河为池。据亿丈之城。临不测之溪以为固。良将劲弩。守要害之处。信臣精卒。陈利兵而谁何。天下已定。始皇之心。自以为关中之固。金城千里。子孙帝王万世之业也。\\
\hspace*{2em}﹁始皇既没。余威振于殊俗。然而陈涉瓮牖绳枢之子。甿隶之人。而迁徙之徒也。材能不及中人。非有仲尼。墨翟之贤。陶朱。猗顿之富也。蹑足行伍之间。俯仰仟佰之中。率罢散之卒。将数百之众。转而攻秦。斩木为兵。揭竿为旗。天下云集响应。赢粮而景从。山东豪俊遂并起而亡秦族矣。\\
\hspace*{2em}﹁且天下非小弱也。雍州之地。殽函之固自若也。陈涉之位。非尊于齐。楚。燕。赵。韩。魏。宋。卫。中山之君也。锄耰棘矜。非铦于勾戟长铩也。适戍之众。非俦于九国之师也。深谋远虑。行军用兵之道。非及乡时之士也。然而成败异变。功业相反也。尝试使山东之国与陈涉度长絜大。比权量力。则不可同年而语矣。然而秦以区区之地。致万乘之权。抑八州而朝同列。百有余年矣。然后以六合为家。殽函为宫。一夫作难而七庙堕。身死人手。为天下笑者。何也。仁义不施。而攻守之势异也。﹂\\

% ===== 索隐述赞（司马贞《史记索隐》）· 统一样式 =====
% 题头「索隐述赞曰」五字并入反白框；正文每四句一列，定宽盒 \hfil 居列中；
% \fontsize{24}{41} 列距对齐全书界行（121+41k）；\clearpage 使述赞独占一页。
\clearpage
\shuzanhead{索隐述赞曰}\linebreak
{\shuzanbody\fontsize{24}{41}\selectfont
\hbox to \hsize{\hfil 天下匈匈。海内乏主。掎鹿争捷。瞻乌爰处。\hfil}\linebreak
\hbox to \hsize{\hfil 陈胜首事。厥号张楚。鬼怪是凭。鸿鹄自许。\hfil}\linebreak
\hbox to \hsize{\hfil 葛婴东下。周文西拒。始亲硃房。又任胡武。\hfil}\linebreak
\hbox to \hsize{\hfil 夥颐见杀。腹心不与。庄贾何人。反噬城父。\hfil}%
\par}
\end{document}
